\baselineskip12pt plus.2pt minus.2pt

\spis{Co nazwy pierwiastków mówią o~stosunku uczonych do nauki}
{Tomasz Ratajczyk}

\wtyt{Co nazwy pierwiastków mówią\\ o~stosunku uczonych do nauki}

\mtyt{Próba systematycznego ujęcia}

\waut{Tomasz RATAJCZYK*}
\aafil{Doktorant, Wydział Chemii Uniwersytetu Warszawskiego,
Szkoła Doktorska Nauk Ścisłych i~Przyrodniczych}

%\emph{Autor poprosił, żeby wytłuszczone były tylko pierwsze wystąpienia każdego pierwiastka (łacznie 118 wytłuszczeń), zaś kursywą tylko źródłosłowy niepolskie.}

Choć minęło zaledwie osiem lat od przyjęcia przez IUPAC ostatnich
czterech nazw dla superciężkich pierwiastków, elegancki wygląd pełnego
siódmego okresu i~brak dalszych doniesień o~sztucznym wytworzeniu
jeszcze cięższych jąder powodują, że zaczynamy myśleć o~układzie
okresowym w~obecnej formie jako o~czymś jak najbardziej kompletnym.
Nawet jeśli tak ostatecznie nie będzie, to poczucie kompletności skłania
do podsumowań, do podjęcia próby pewnej systematyki i~wyciągnięcia z~niej wniosków, podobnie jak to czyniono w~przypadku innych kompletnych
systemów, np. układu czterech żywiołów Empedoklesa, pięciu brył
platońskich albo modelu standardowego cząstek elementarnych. Tym, co
wyróżnia układ okresowy spośród nich, jest dość duży zbiór elementów, a~przede wszystkim niezwykle długi proces nadawania im nazw, trwający od
prehistorii do czasów najnowszych.

%\begin{dwieszpalty}
\ptyt{{Początki historii}}

Najstarsze nazwy obejmują pierwiastki, które były znane najdawniej, i~znaczą dokładnie to, co znaczą. \textbf{Węgiel} jaki jest, każdy widzi,
tak też i~\textbf{siarka}. Korzenie tych słów sięgają języka
praindoeuropejskiego, więc są wielokrotnie starsze niż samo pojęcie
pierwiastka chemicznego. Utożsamienie tych substancji z~pierwiastkami
nic w~nich nie zmieniło w~sensie fizykalnym. Podobnie siedem znanych w~starożytności metali, kolejny kompletny układ, sprzężony z~siedmioma
znanymi wówczas ciałami niebieskimi i~siedmioma wspólnymi dla wielu
kultur bóstwami: \textbf{srebro}--Księżyc, \textbf{rtęć}--Merkury,
\textbf{miedź}--Wenus, \textbf{złoto}--Słońce, \textbf{żelazo}--Mars,
\textbf{cyna}--Jowisz, \textbf{ołów}--Saturn. Ten system wydawał się do
tego stopnia kompletny, że jeszcze w~nowożytności nowo odkryte metale
przypisywano do któregoś z~tych siedmiu. Na przykład \textbf{molibden}
to po grecku ołów (\emph{molybdos}), do którego podobny wizualnie jest
molibdenit (MoS$_2$). \textbf{Platyna} -- to z~hiszpańskiego srebro
(\emph{plata}), i~to pogardliwie zdrobnione, bowiem dawni Hiszpanie,
podbijając Amerykę Środkową, uznali tamtejsze platynowe skarby za gorszy
sort tego metalu. Po drodze pojawiło się jeszcze kilka nazw o~dość
mglistej etymologii: dwie perskie, \textbf{arsen} i~\textbf{cyrkon},
oraz dwie staroniemieckie, \textbf{cynk} i~\textbf{bizmut}. Jednak
dopiero Robert Boyle, dając pierwszą ścisłą definicję pierwiastka, dał
zarazem asumpt do dalszych, dużo bardziej świadomych odkryć\ldots{}
które trzeba było zacząć bardziej świadomie nazywać.

\ptyt{{Coś musiało być w~tych minerałach}}

Kolejną kategorię stanowią pierwiastki nazwane od substancji, z~których
je pozyskano, nieposiadających głębszych źródłosłowów. Są to głównie
pierwiastki lekkie i~wcześnie odkryte; niektóre z~nich Antoine Lavoisier
ponazywał po francusku jeszcze przed wydzieleniem ich w~postaci czystej,
słusznie przewidując, że w~konkretnych minerałach tkwią konkretne,
odrębne pierwiastki. Po polsku mamy \textbf{sód} od sody, \textbf{potas}
od potażu (ang. \emph{potash}, czyli popiół roślinny), \textbf{lit} od
gr. \emph{lithos} (kamień), bo w~przeciwieństwie do potasu, otrzymano go
z~ziemi, z~jakiegoś kamienia. Dalej: \textbf{beryl} od minerału berylu,
\textbf{wapń} od wapna, \textbf{magnez} od magnezji. Od magnezji, albo
raczej starożytnego miasta Magnezji w~Grecji, pochodzi nie tylko magnez,
ale i~\textbf{mangan}, jednak nazywając te pierwiastki, o~tym miejscu
nikt już raczej nie pamiętał. Jeszcze dalej: \textbf{bor} od boraksu,
\textbf{glin} od glinki (łac. \emph{alumina}), \textbf{krzem} od
krzemionki (łac. \emph{silica}). Wiele z~tych polskich nazw zawdzięczamy
Jędrzejowi Śniadeckiemu, który za pokolenie starszym Lavoisierem porwał
się na wprowadzenie systematycznych polskich nazw, wzorowanych na
francuskich. Jeszcze jedną z~nich jest \textbf{wodór}, który przecież
wziął się stąd, że tworzy wodę.\eject
\begin{dwieszpalty}
\ptyt{{Właściwości pierwiastka\ldots{}}}

Dotychczas omówione nazwy pierwiastków są w~pewnym sensie podstawowe,
ponieważ za ich główną właściwość wzięto to, że są tym, czym są. W~następnej, najobszerniejszej grupie znajdą się pierwiastki, które
zostały nazwane od jakiejś konkretniejszej właściwości, fizycznej,
chemicznej lub symbolicznej, którą uznano za najbardziej
charakterystyczną. Na przykład \textbf{fosforowi} zdarza się świecić,
czyli „nieść światło'', podobnie jak semafor „niesie sygnały''. Za
najbardziej charakterystyczne dla \textbf{azotu} przyjęto to, że nie
podtrzymuje życia (gr.~\emph{azotikos}), przynajmniej w~Polsce, Rosji
czy Francji, bo już Niemcy za istotniejsze uznali niepodtrzymywanie
spalania (niem. \emph{Stickstoff}), a~anglofoni -- występowanie tego
pierwiastka w~saletrze (łac. \emph{nitrum}, stąd \emph{nitrogen}).
Spalanie (bardziej poetycko \emph{tlenie}) podtrzymuje za to
\textbf{tlen}, w~języku polskim nazwany wyjątkowo trafnie przez Jana
Oczapowskiego w~połowie XIX wieku, zamiast już wtedy nieaktualnego
kwasorodu, który był kalką z~fr. \emph{oxygène}. Nieliczne pierwiastki
mają w~czystej postaci nietypowe kolory, co zgrabnie wykorzystano w~nazewnictwie: \textbf{chlor} to po grecku zielony (\emph{chloros}), a~\textbf{jod} -- fioletowy (\emph{iodes}). \textbf{Brom} wprawdzie mógłby
być brunatny, ale za bardziej charakterystyczne uznano, że cuchnie (gr.~%
\emph{bromos}), pod czym zresztą ciężko się nie podpisać. Od okropnej
woni nazwano aż dwa pierwiastki, tyle że w~przypadku drugiego z~nich,
\textbf{osmu}, tak naprawdę przykry zapach (gr.~\emph{osme}) ma
powstający na powierzchni lotny tlenek.

\ptyt{{\ldots albo jego związków}}

Cechy wybranych związków również bywały źródłami nazw ich składników.
Niektóre jeszcze poprzez znane od dawna minerały, choćby \textbf{fluor}
od fluorytu (łac. \emph{fluo}, płynę), dodawanego w~hutach jako topnik.
Z~drugiej strony, wolframit w~wytapianiu rud przeszkadzał, ponieważ
prowadził do uciążliwego („wilczego'', niem. \emph{Wolf}) powstawania
piany (niem. \emph{Rahm}), stąd dziś metal \textbf{wolfram}. Górnikom
dawały się we znaki także inne złośliwe moce, takie jak potworki Kobold
(stąd \textbf{kobalt}) i~Nickel (stąd \textbf{nikiel}), utrudniające
prace, odpowiednio, z~żelazem i~z~miedzią. Później, ze względu na
właściwości ich związków, nazwano też \textbf{chrom} kolorowym (gr.~%
\emph{chroma}, kolor), \textbf{rod} różowym (gr.~\emph{rhodon}, róża),
\textbf{iryd} tęczowym (łac. \emph{iris}, tęcza), a~\textbf{bar} ciężkim
(gr.~\emph{barys}). Podobnie \textbf{prazeodym}, który niegdyś
traktowano razem z~\textbf{neodymem} jako jeden pierwiastek, didym (gr.~%
\emph{didymos}, bliźniak), nazwany tak, gdyż był bliźniaczo podobny do
lantanu i~ceru. Ironia losu, że jakiś czas później okazało się, że didym
sam w~sobie stanowi mieszaninę dwóch różnych pierwiastków. Gdy już się
udało rozdzielić tych bliźniaków, jeden z~nich okazał się tworzyć
związki o~barwie zielonej (gr.~\emph{prasios}, zielony), a~drugi po
prostu nazwano nowym (gr.~\emph{neos}).

\ptyt{{Małe przełomy}}

Niejednokrotnie seria podobnych odkryć następowała po powstaniu
przełomowego wynalazku. Wspomniane sód, wapń itp. wydzielono jeden po
drugim dzięki elektrolizie. Inne z~kolei dzięki spektroskopii.
Porównując widma emisyjne, odkryto czerwony prążek \textbf{rubidu}
(dzielącego łac. źródłosłów \emph{rubeus} z~czerwonymi rubinami),
niebieskie prążki \textbf{indu} (\emph{indygo}) czy \textbf{cezu} (łac.
\emph{caesius}, błękitny), zielony prążek \textbf{talu} (gr.
\emph{thalos}, zielona gałązka). Tak odkryto także \textbf{hel}, nazwany
od Słońca (gr. \emph{Helios}), w~którego widmie znaleziono prążki
pochodzące od tego pierwiastka. Pozostałe gazy szlachetne znaleziono
równie szybko dzięki rozwojowi kriogeniki. Ponieważ każdy z~nich wydawał
się dość podobny, a~jednocześnie inertny i~trudny do wyizolowania, dano
im z~greki nieco symboliczne imiona: nowy \textbf{neon}, leniwy
\textbf{argon}, ukryty \textbf{krypton} i~obcy \textbf{ksenon}. Jeszcze
kilka innych pierwiastków nazwano od trudności w~otrzymaniu:
\textbf{lantan} i~\textbf{dysproz} na pewno (gr. \emph{lanthano},
ukrywam się, \emph{dysprositos}, uciążliwy), a~w~pewnym sensie także
\textbf{technet} (gr. \emph{technetos}, sztuczny) i~\textbf{antymon}
(gr. \emph{antimonos}, niesamodzielny). Wreszcie po odkryciu
promieniotwórczości pojawiło się kilka pierwiastków nazwanych od niej:
\textbf{rad} (a~po nim \textbf{radon}) i~\textbf{aktyn} (a~przed nim
\textbf{protaktyn}), no i~niestabilny \textbf{astat} (gr.
\emph{astatos)}.

\ptyt{{Odniesienia do mitologii}}

A~co, jeżeli nie udało się wyłonić jednej istotnej właściwości? Z~początku czasami sięgano po symbolikę ugruntowaną w~mitologii, głównie
grecko-rzymskiej. Z~tej zaczerpnięto choćby \textbf{tytan} (intuicyjnie
brzmi jednak jak coś wytrzymałego, prawda?), występujące często razem
\textbf{niob} i~\textbf{tantal} (Niobe była córką Tantala),
\textbf{kadm} pośrednio od króla Kadmosa, \textbf{pallad} od Pallas
Ateny. \textbf{Promet} nazwano od Prometeusza, który wykradł z~Olimpu
ogień -- pod ogień należy tu wstawić energię jądrową. Kojarzący się z~nią \textbf{uran} akurat został nazwany jeszcze w~XVIII wieku po
odkrytej przez Williama Herschela niewiele wcześniej planecie, ale już odkryte
podczas II wojny światowej \textbf{neptun} i~zwłaszcza \textbf{pluton}
-- oprócz mitologii pobrzmiewają też echami potężnego zniszczenia.
Cztery inne nazwy pochodzą z~mitologii germańskiej: \textbf{tor} od
równie silnego Thora, \textbf{wanad} od bogini Vanadis% (jedno z~imion Freyi)
, a~także wymienione już kobalt i~nikiel od duchów bardziej
wrednych niż potężnych. Mniej destrukcyjnie kojarzą się za to
\textbf{cer}, zaraz po planetoidzie nazwany tak od rzymskiej bogini
Ceres, a~także \textbf{selen} od gr. \emph{Selene} (Księżyc) i~\textbf{tellur} od łac. \emph{Tellus} (Ziemia). Podobnie jak w~przypadku
klasycznych siedmiu metali, i~tutaj czasem trudno określić, czy
więcej inspiracji dało ciało niebieskie, czy odpowiadająca mu postać z~mitologii.

\ptyt{{Gdzie są te pierwiastki?}}

Przedostatnia duża grupa nazw obejmuje te pochodzące od nazw miejsc
(toponimów). Wiele  pierwiastków z~tej grupy odkryto względnie niedawno i~w~większości motywacją do nadania nazwy była chęć uhonorowania miejsca
odkrycia pierwiastka albo pochodzenia odkrywcy. W~dziewiętnastym wieku
dominowała %tu
w~tym Europa, do której naturalnie dołączyły później USA.
Skandynawski jest oczywiście \textbf{skand}, ale też cztery
\end{dwieszpalty}\eject pierwiastki
z~kopalni Ytterby w~Szwecji, czyli \textbf{itr}, \textbf{iterb},
\textbf{terb} i~\textbf{erb}, dalej \textbf{tul} (gr. \emph{Thule},
półlegendarna nazwa nordyckiej wyspy), \textbf{holm} od Sztokholmu i~\textbf{hafn} od Kopenhagi. \textbf{Stront} nazwano od szkockiej wsi
Strontian. \textbf{German} od Germanii, czyli Niemiec, w~których
znajdują się też kraj związkowy Hesja (\textbf{has}), miasto Darmstadt
(\textbf{darmsztadt}) oraz rzeka Ren (\textbf{ren}). Są i~trzy
francuskie pierwiastki: \textbf{frans}, \textbf{gal} od Galii i~\textbf{lutet} od Lutecji, osady, wokół której zbudowano Paryż.
\textbf{Polon}, choć też odkryty we Francji, nazwała Maria Skłodowska,
więc Polakom również udało się zapisać w~układzie okresowym. Śniadecki
sto lat wcześniej także był na tropie
% pierwiastka, ale w~końcu mu to nie wyszło, a~pierwiastek
pewnego pierwiastka, ale jego odkrycia nie udało się potwierdzić, a~później pierwiastek ten
nazwano \textbf{rutenem} -- od Rusi. 
Oprócz rutenu
jeszcze dwa pierwiastki mają eponimy rosyjskie: \textbf{dubn} od
instytutu w~Dubnej i~\textbf{moskow} od Moskwy, blisko której znajduje
się ten instytut. Po odnotowaniu przy tym oczywistego \textbf{europu}
należy zwrócić uwagę na Amerykę (\textbf{ameryk}), gdzie w~stanie
Kalifornia (\textbf{kaliforn}) pracują instytuty Berkeley
(\textbf{berkel}) i~Livermore (\textbf{liwermor}); jest także stan
Tennessee, znany nie tylko z~burbonu, ale i~z~pierwiastka o~nazwie
\textbf{tenes}. Japończykom udało się nazwać tylko jeden --
\textbf{nihon}.


\ptyt{{I~ty możesz zostać pierwiastkiem}}

Wraz z~rozwojem metod otrzymywania superciężkich jąder nowo wytworzone
pierwiastki zaczęły znajdować się na granicy odkrycia i~wynalazku. Czasy
życia pojedynczych atomów, mierzone w~milisekundach, nie pozwalają i~być
może nigdy nie pozwolą na zbadanie makroskopowych właściwości
transkiurowców. Zatem w~nazwach najcięższych pierwiastków obok toponimów
przeważają antroponimy, których prawie nie ma wśród pierwiastków
stabilnych
(wyjątkami są \textbf{samar} od Wasilija Samarskiego przez minerał samarskit oraz \textbf{gadolin} od Johana Gadolina przez minerał gadolinit).
%(wyjątkami są \textbf{samar} i~\textbf{gadolin}, których nazwy pochodzą od nazw minerałów gadolinitu i~samarskitu nazwanych z~kolei na cześć badających je uczonych Wasilija Samarskiego i~Johana Gadolina
%s{amar} od Wasilija Samarskiego przez minerał samarskit oraz {gadolin} od Johana Gadolina przez minerał gadolinit).%
Poza nimi pozostałych 14 nazwisk znalazło się w~ostatnim
okresie układu. Można je wymienić ciągiem: \textbf{kiur} (Maria
Skłodowska-Curie), \textbf{einstein} (Albert Einstein), \textbf{ferm}
(Enrico Fermi), \textbf{mendelew} (Dmitrij Mendelejew), \textbf{nobel}
(Alfred Nobel), \textbf{lorens} (Ernest Lawrence), \textbf{rutherford}
(Ernest Rutherford), \textbf{seaborg} (Glenn Seaborg), \textbf{bohr}
(Niels Bohr), \textbf{meitner} (Lise Meitner), \textbf{roentgen}
(Wilhelm Roentgen), \textbf{kopernik} (Nicolaus Copernicus),
\textbf{flerow} (Gieorgij Florow), \textbf{oganeson} (Jurij Oganiesian).
Na wyróżnienie zasługuje kopernik, nazwany po uczonym względnie bardzo dawnym, a przeciwnie seaborg i oganeson, nazwane po tych naukowcach jeszcze za ich życia.
% Na wyróżnienie zasługują: kopernik, nazwany po uczonym względnie bardzo dawnym, i~przeciwnie -- seaborg i~oganeson, nazwane po naukowcach jeszcze za ich życia.

\ptyt{{Powrót do początków}}

Porównując wszystkie zebrane wyżej nazwy pierwiastków, można wskazać
kilka głównych tendencji. Najstarsze nazwy wykształcały się w~sposób
naturalny razem z~całym językiem. Jeśli pierwiastek był znany jako taki,
oznaczały ten pierwiastek, jeśli znany był jego związek, oznaczały ten
związek. Stulecia XVIII, XIX i~początek XX przyniosły wiele neologizmów
świadomie utworzonych przez badaczy, zazwyczaj zaczerpniętych z~greki
lub łaciny, odnoszących się do właściwości opisywanej substancji, a~czasami do mitologii. Kiedy pierwiastków znano coraz więcej i~okazało
się, że coraz więcej z~nich jest do siebie podobnych, coraz częściej
nazywano je od różnych miejsc. Pod koniec XIX wieku zaczęło też zanikać
wykształcenie klasyczne, więc naukowcom źródłosłowy greckie czy
łacińskie po prostu stopniowo przestały przychodzić do głowy. Wreszcie
do nazw odmiejscowych dołączyły nazwiska zasłużonych badaczy, mniej lub
bardziej związanych z~rozszerzaniem układu okresowego. Superciężkie
pierwiastki nie mają właściwości dających się łatwo zaobserwować
„szkiełkiem i~okiem'', wykorzystanie nazwisk wydaje się naturalną tego
konsekwencją.

W~pewnym sensie wróciliśmy więc do punktu wyjścia. „Ciężki'' bar czy
„niebieski'' ind to nazwy o~bardzo wdzięcznym pochodzeniu, które jednak
podpowiada nam jakąś jedną właściwość pierwiastka. Niekoniecznie
najważniejszą i~niekoniecznie najtrafniejszą, jak w~przypadku
tlenu/kwasorodu. Liwermor czy tenes też niby od czegoś pochodzą, ale
właściwie w~kontekście tego, czym są te pierwiastki, to nic nie znaczy.
Tenes jest tym, czym jest tenes. Węgiel jest tym, czym jest węgiel.
\marg[-160]{\parskip2pt

  \textbf{Literatura}

I. Eichstaedt, \emph{Księga pierwiastków}, Wiedza Powszechna, Warszawa
1973.

A. Lavoisier, \emph{Traktat podstawowy chemii}, tł. R. Mierzecki,
Analecta X, 2001, 1(19), 7--122.

H. Lichocka, \emph{Chemia Jędrzeja Śniadeckiego}, Opuscula Musealia 26,
2019, 21--34.

R. Mierzecki, \emph{Rozwój polskiej terminologii chemicznej}, Zakład
Narodowy im. Ossolińskich, Wrocław 1988.

A. Sulich, \emph{Nomenklatura pierwiastków chemicznych w~językach
świata}, Acta Philologica 42, 2012, 5--13.

A. Sulich, \emph{Odwzorowanie wiedzy w~nazwach pierwiastków
chemicznych}, Studia Semiotyczne XXVII, 2010, 93--133.}

\eject
%\end{dwieszpalty}
\normalsize